\documentclass[../../main.tex]{subfiles}
\begin{document}

{\bf [解答]} $f(x)$を$f$などと略記する．題意の条件から
\begin{numcases}
  {}
  u(x)=\int_{0}^{x}h(t)f(t)dt+h\int_{x}^{1}f(t)dt \label{eq:1}\\
  u''(x)=f(x) \label{eq:2}\\
  u(0)=0\label{eq:3}
\end{numcases}
まず，\cref{eq:1}で$x=0$として,\cref{eq:3}から
\begin{align}
  0=h(0)\int_{0}^{1}fdt
\end{align}
ここで$f$は任意の一次関数だから$\int_{0}^{1}fdt\equiv0$とはならないので(例えば$f=x$)
\begin{align}
  h(0)=0\label{eq:4}
\end{align}
である．

\cref{eq:1}の両辺を$x$で微分する．
\begin{align}
  u'&=hf+h'\int_x^1fdt-hf  \\
  &=h'\int_x^1fdt
\end{align}
さらに微分して
\begin{align}
  u''=h''\int_x^1fdt-h'f
\end{align}
\cref{eq:2}を代入して
\begin{align}
  \{1+h'\}f=h''\int_x^1fdt\label{eq:5}
\end{align}
$f(x)$は任意の一次関数で，$a \ (\neq 0),b\in\mathbb{R}$として，\cref{eq:5}に$f(x)=ax+b$を代入したとき両辺$a$で割ることにより，$f(x)=x+b$のみ考えれば良いことがわかる．実際に代入して
\begin{align}
  &\{1+h'\}(x+b)=\left[\dfrac{1}{2}t^2+bt\right]_x^1h'' \\
  =&\{\dfrac{1}{2}(1-x^2)+b(1-x)\}h''
\end{align}
これが$b$についての恒等式だから係数比較して，
\begin{subnumcases}
  {}
  1+h'-(1-x)h''=0 \label{eq:6}\\
  (1+h')x-\dfrac{1}{2}h''(1-x)(1+x)=0 \label{eq:7}
\end{subnumcases}
となる．\cref{eq:6}を\cref{eq:7}に代入して
\begin{align}
  &(1+h')x-\dfrac{1}{2}(1+h')(1+x)=0 \\
  \Longleftrightarrow
  &(1+h')(x-1)=0
\end{align}
これが$x$についての恒等式だから，$h'=-1$が従う．\cref{eq:4}と合わせて
\begin{align}
  h(x)=-x
\end{align}
である．これが\cref{eq:1,eq:2,eq:3}を満たすことは容易に確かめられる．ゆえに求める関数は
\begin{align}
  h(x)=-x
\end{align}
である．$\cdots$(答)

\newpage
\end{document}
